\section{Introduction}\label{sec:intro}

The $n$-queens problem asks for a placement of $n$ queens on an $n\times n$ board such that no two queens
attack each other. A solution
has exactly one queen in each row and each column and at most one on each diagonal in either direction. We
denote the number of
solutions by $Q(n)$ (sequence A000170 in~\cite{OEIS}); Bell and Stevens~\cite{BellStevens} survey the known
results on the
problem.

Consider a board one step from completion. Once $n-1$ nonattacking queens have been placed, exactly one row
and one column remain
empty. Their intersection is the only square where the last queen could be placed without moving the others.
The placement can be
completed if and only if both diagonals through that square are free. If the placement of the $n-1$ queens
is chosen uniformly at
random, what is the probability of this event as the board grows?

We use small boards as examples. In the figures, rows and columns are numbered from zero, with the row
coordinate given first.
Figure~\ref{fig:ejemplos} shows a solution on the $4\times4$ board and another on the $5\times5$ board.

\begin{figure}[htbp]
\centering
\begin{tikzpicture}[tablero A]
 \begin{scope}
  \cuadriculaA{4}\trazarRejillaA{4}
  \reinaA{0}{1}\reinaA{1}{3}\reinaA{2}{0}\reinaA{3}{2}
  \panelA{2,5.34}{(a) A $4\times4$ solution}
 \end{scope}
 \begin{scope}[shift={(7,0)}]
  \cuadriculaA{5}\trazarRejillaA{5}
  \reinaA{0}{1}\reinaA{1}{3}\reinaA{2}{0}\reinaA{3}{2}
  \aroA{4}{4}\reinaA[colorAro]{4}{4}
  \panelA{2.5,5.34}{(b) A $5\times5$ solution}
 \end{scope}
\end{tikzpicture}
\caption{Two small solutions. Each crown represents a queen. The ring marks the queen at $(4,4)$ used in the examples that follow.}
\label{fig:ejemplos}
\end{figure}


Throughout the paper we fix the board size and group placements by the number of missing queens. This
explains the word \emph{layer} in
the title. For example, removing the queen $(4,4)$ in Figure~\ref{fig:ejemplos}(b) gives a placement in
layer $h=1$ of the $5\times5$ board.
Putting it back takes us to layer $h=0$.

The same question makes sense when $h$ queens are missing. We then say that the placement has $h$
\emph{holes}. We call the set
$\Om_h(n)$ of legal placements of $n-h$ queens \emph{layer $h$}. Let $N_h(n)=|\Om_h(n)|$ and let $\mu_h$ be
the uniform law on this
layer whenever it contains at least one placement. Take a placement with law $\mu_h$. Its $h$ empty rows and
$h$ empty columns
intersect in $h^2$ \emph{empty crossings}. By the positions of the holes we shall mean the positions of
these crossings. Choose one of them
uniformly at random. We say it is \emph{available} if a queen can be added there, which is equivalent to
both of its diagonals being
free. For $h>1$, this experiment concerns the next insertion, even though some queens will still remain to
be placed.

The probability of a legal insertion can be expressed in terms of the sizes of two consecutive layers. Each
placement in layer $h-1$
has $n-h+1$ queens that we can remove individually to return to layer $h$. Conversely, each placement in
layer $h$ offers $h^2$
crossings at which to try an insertion. Counting legal insertions in both directions
(Lemma~\ref{lem:interpretacion}) shows that the
mean fraction of available crossings under $\mu_h$ is the ratio \begin{equation}\label{eq:fh}
 f_h(n)=\frac{(n-h+1)\,N_{h-1}(n)}{h^2\,N_h(n)} ,
\end{equation} which we take as the definition of $f_h(n)$ for $1\le h\le n$. Thus $f_1(n)$ answers the
question about the last queen
exactly. We prove that for each fixed $h$, $f_h(n)$ tends to the same positive constant as the board grows.

To see what is being averaged, return to Figure~\ref{fig:ejemplos}(b) and remove the queens $(1,3)$ and
$(4,4)$. With three queens left
on the $5\times5$ board, there are two holes and four crossings (Figure~\ref{fig:cruces-vacios}). The empty
rows are $1$ and $4$, and the empty
columns are $3$ and $4$. The queen at $(3,2)$ attacks $(1,4)$ along the antidiagonal with sum $5$ and
$(4,3)$ along the diagonal with
difference $1$. Each of the other two crossings can receive a queen.

\begin{figure}[htbp]
\centering
\begin{tikzpicture}[tablero A]
 \cuadriculaA{5}
 \foreach \rA in {1,4}{\fill[colorAro!7] (0,\rA) rectangle ++(5,1);}
 \foreach \cA in {3,4}{\fill[colorAro!7] (\cA,0) rectangle ++(1,5);}
 \foreach \rA in {1,4}{\foreach \cA in {3,4}{\fill[colorFondo] (\cA,\rA) rectangle ++(1,1);}}
 \trazarRejillaA{5}
 \ataqueA{3}{2}{1}{4}\ataqueA{3}{2}{4}{3}
 \reinaA{0}{1}\reinaA{2}{0}\reinaA{3}{2}
 \aroA{1}{3}\aroA{4}{4}\aspaA{1}{4}\aspaA{4}{3}
 \begin{scope}[shift={(6.0,.23)}]
  \node[anchor=west,font=\small\sffamily\bfseries] at (0,.25) {The four crossings};
  \aroA{1}{0}
  \node[anchor=west,font=\small] at (1.06,1.5) {Available: $(1,3)$ and $(4,4)$};
  \aspaA{2}{0}
  \node[anchor=west,font=\small] at (1.06,2.5) {Attacked: $(1,4)$ and $(4,3)$};
  \node[anchor=west,font=\small] at (.1,3.95) {$X_2(P)=2$ of $2^2=4$ crossings};
 \end{scope}
\end{tikzpicture}
\caption{A placement in layer $\Omega_2(5)$. The circles and crosses mark the four intersections of the
empty rows and columns. The dashed lines show why two of these squares are under attack.}
\label{fig:cruces-vacios}
\end{figure}


We write $X_h(P)$ for the number of available crossings in a placement $P$. In the example, $X_2(P)=2$, so a
uniformly chosen crossing
on this board can receive a queen with probability $2/4$. To obtain $f_2(5)$, we must average this fraction
over all legal placements
of three queens on the $5\times5$ board.

\subsection{Main results}

To identify the constant, we write $x=h/n$ for the fraction of holes. As $n$ grows with $h$ fixed, $x$ tends
to zero. This leads us to
study the full-board endpoint of the variational profile, namely the density $\nu_x$ on $[0,1]^2$ with mass
$1-x$ that maximizes the
entropy functional of Section~\ref{subsec:variacional}. For $x\in[0,1/2]$ it has the form \[
 \nu_x(r,c)=b_x\,r_{R,x}(r)\,r_{C,x}(c)\,u_{D,x}(r-c)\,u_{A,x}(r+c),
 \qquad \int_0^1 r_{R,x}=\int_0^1 r_{C,x}=1,
\] where $r_{R,x}$ and $r_{C,x}$ describe the rescaled profiles of the empty rows and columns, $u_{D,x}$ and
$u_{A,x}$ describe the
\emph{vacancies} of the diagonals and antidiagonals (the density of those that remain empty), and $b_x>0$ is
the normalization
constant. The properties of the profile are proved in~\cite{Perfil}. We define \begin{equation}\label{eq:F}
 F(x)=\frac{1-x}{b_x}=\int_{[0,1]^2}r_{R,x}(r)\,r_{C,x}(c)\,u_{D,x}(r-c)\,u_{A,x}(r+c)\,dr\,dc .
\end{equation} Heuristically, $F(x)$ is the probability that both diagonals through a typical empty crossing
are free when there are
$xn$ holes. The position of the crossing has law $r_{R,x}r_{C,x}$, and each diagonal is empty with density
$u_{D,x}$ or $u_{A,x}$
according to its direction. For the final layers the relevant constant is $F(0)=1/b_0$.

\begin{teorema}[Law of the final layers]\label{thm:limite} With the profile of Theorem~\ref{thm:perfil},
taken
from~\cite{Perfil}, for each fixed $h\ge1$, \[
 \lim_{n\to\infty}f_h(n)=F(0)=\frac1{b_0}\in(0,\infty).
\] In particular, $\max_{1\le h\le K}|f_h(n)-F(0)|\to0$ for each fixed $K$. \end{teorema}

Each factor $f_j(n)$ relates two consecutive layers. Multiplying these relations from $j=1$ to $j=h$ brings
the number of complete
solutions into the denominator.

\begin{corolario}[Ratio of layer sizes]\label{cor:cociente} For each fixed $h\ge1$, \[
 \frac{N_h(n)}{Q(n)}=\frac{(n)_h}{(h!)^2\prod_{j=1}^hf_j(n)}
 \sim\frac{n^h}{(h!)^2F(0)^h},\qquad (n)_h=n(n-1)\cdots(n-h+1).
\] \end{corolario}

Passing from layer $h$ to layer $h-1$ changes the fraction of holes from $h/n$ to $(h-1)/n$. The next
corollary uses as reference
the value of $F$ at the midpoint, $(h-\tfrac12)/n$.

\begin{corolario}[Fixed layers contribute no residual constant]\label{cor:capas} For each fixed $K$, \[
 \sum_{h=1}^K\log\frac{f_h(n)}{F\bigl((h-\tfrac12)/n\bigr)}\longrightarrow0,
\] and the same holds with $F(h/n)$ in place of $F((h-\tfrac12)/n)$. \end{corolario}

The corollary is Theorem~\ref{thm:limite} written with logarithms: each summand tends to zero because
$f_h(n)\to F(0)$ and $F$ is
continuous at $0$. Hence it does not single out the midpoint, and the conclusion holds equally with
$F(x_{n,h})$ for any
$x_{n,h}\to0$.

The number $K$ is fixed before the board size tends to infinity. Identifying the constant term of
$\log Q(n)$ would also require
control over layers whose index grows with $n$. For the final two layers we also give a direct proof of an
explicit lower bound on
$\liminf_nf_h(n)$ that does not require concentration of the reverse count (Section~\ref{sec:inversa}).

\begin{teorema}[Lower bound with an explicit constant]\label{thm:positiva} Let $\kappa\in(0,1]$ be the lower
bound on the diagonal
vacancies $u_{D,x}$ and $u_{A,x}$ of the profile of Theorem~\ref{thm:perfil}, taken from~\cite{Perfil}, and
let
$a_0=(\kappa/2)^4/16$. Then \[
 \liminf_{n\to\infty}f_h(n)\ge\frac{a_0}{2}>0,\qquad h=1,2.
\] In particular, there exist $c>0$ and $n_0$ such that $f_1(n),f_2(n)\ge c$ for all $n\ge n_0$.
\end{teorema}
